\chapter{Milestone log}
\label{ch:milestones}

\begin{quote}\itshape
The project's dated gate-by-gate record, kept as written at each milestone (entries are
historical documents; later chapters supersede their numbers where noted).
\end{quote}


\begin{table}[htbp]
  \centering\small
  \caption{Milestones reached (updated at every gate).}
  \label{tab:milestone-log}
  \begin{tabular}{lp{0.11\textwidth}p{0.62\textwidth}}
    \toprule
    Milestone & Date & Evidence \\
    \midrule
    H0 & 2026-07-03 & Reproducible environment (uv, Python 3.11, pyCycle 4.4 with
      \code{numpy<2}); a complete example cycle converging locally and in CI; repository
      skeleton and smoke test. \\
    H1 (partial) & 2026-07-03 & WP1.1 and WP1.2 complete: design point inside all five
      acceptance windows (\cref{tab:design}); SLS anchors calibrated against TCDS + EEDB with
      all three gated anchors in tolerance (\cref{tab:anchors}); calibration contract with
      verified values. \\
    H1 (WP1.3--1.4) & 2026-07-03 & Baseline decks over the flight envelope with leave-one-out
      quality report (\cref{tab:deck-loo}); health-parameter injection validated by the
      off-design/design self-consistency check; ICM at both snapshot conditions with SVD
      observability analysis --- minimum cockpit signature angle 1.3$^{\circ}$, seven
      confusable pairs (\cref{tab:confusable}); tabular-vs-CEA cross-check within 0.7\,\%
      (\cref{tab:cea}); twelve automated tests green. \\
    \textbf{H1 closed} & 2026-07-03 & Correction exponents fitted from the model land on
      textbook values (\cref{tab:corrected}); corrected-space baseline collapses LOO tails by
      an order of magnitude; ICM extended to a six-point envelope grid with stable
      observability (\cref{tab:icm-grid}); computations parallelized per norm N1 (decks
      $2.9\times$, ICM $6\times$ faster); sixteen automated tests green. \\
    \textbf{H2 closed} & 2026-07-03 & SynCFM56 v1.0 frozen (\cref{tab:fleet}): 100 engines,
      1.58M snapshot rows, deterministic regeneration. Nonlinearity at the noise floor
      (\cref{tab:nonlin}); difficulty gate failed at 81.6\,\% with chronic labels, passed at
      54.0\,\% after the acute-episode design change (\cref{sec:audit-difficulty}); realism
      checks all positive; onset bug found and fixed by the audits; 25 automated tests green. \\
    \textbf{H3 closed} & 2026-07-03 & Traditional pipeline end-to-end on the test fleet
      (\cref{tab:trad}): detection recall 0.5 at 98-cycle median delay, isolation 0\,\% on
      the confusable subset, smearing index 0.85, RUL optimism funnel measured
      (\cref{fig:rul-funnel}); two detector-design lessons recorded; all defaults registered
      in \cref{ch:decisions} and exposed to the F5 budget; 31 automated tests green. \\
    F4 addendum & 2026-07-04 & PCS machinery implemented with a
      recorded null result on the weak v1.1 classifier (\cref{sec:pcs}); Mahalanobis
      detector validation sweep found a structural, not threshold, gap
      (\cref{sec:det-sweep}); H4 explicitly not declared --- resolution moved into the F5
      tuned round. \\
    F4 v0 & 2026-07-04 & AI suite end-to-end on MPS (\cref{tab:ai-vs-trad}): RUL RMSE
      4--6$\times$ below traditional with near-zero bias; conformal coverage
      0.85/0.95/1.00 at 0.9 nominal; naive-AE recall-0 and engine-level sample-starvation
      findings recorded; stacking hybrid \emph{negative} result (\cref{fig:hybrid}) logged
      pre-freeze; 36 automated tests green. \\
    \textbf{F5 verdicts} & 2026-07-04 & Tuned under the frozen 50-trial budget; single
      confirmatory pass; H1 confirmed (detection flips to AI: recall 0.48 vs 0.13, delay
      499 vs 6033 cy), H2 refuted (31\,\% = 31\,\%: the observability wall binds both),
      H3 confirmed (Cliff's $\delta$ 0.89), H4 refuted, H5 confirmed (coverage 0.883,
      half-width 5.9$\times$ tighter). \cref{tab:verdicts}; 37 tests green. \\
    \textbf{H5 closed} & 2026-07-04 & Sim-to-real on C-MAPSS FD001 (disclosed N-CMAPSS
      substitution): RUL ranking preserved, 14.9 vs 42.1 cycles RMSE
      (\cref{sec:sim-to-real}). All five hypotheses adjudicated; evidence reproducible end
      to end. \\
    \textbf{F7} & 2026-07-04 & Operating-point tomography (tag prereg-v2): the
      report schedule as an observability design variable (+79\% separability from one report
      condition); learned MOPA tripled classical stacking (H7.2$'$); conformal isolation sets
      (H7.4$'$); drift-robustness thresholds refuted as frozen (H7.3). \cref{ch:tomography}. \\
    \textbf{F8} & 2026-07-04/05 & Seven limitation-overcoming lines (tags prereg-v3/5/6/7/8):
      L1 differentiable twin, L2 nonlinear fleet v2, L5 architecture-robust RUL, L4 recoverable
      fraction ($R^2$ 0.86); L6 H4 refuted again; L7, L9 partial. \cref{tab:f8-summary}. \\
    \textbf{F10} & 2026-07-05 & Ground-truth-validated identifiability certificate
      (tag prereg-v4): honest per-direction precision ($\rho$ 0.70), CRB shape + conformal
      scale, sensor-acquisition instrument. \cref{ch:breakthrough}. \\
    \textbf{F-econ} & 2026-07-05 & Cautious impact estimate for a 100-engine operator
      (median \$3\,M/yr, wide interval, honest downside). \cref{ch:economics}. \\
    \textbf{Audit} & 2026-07-05 & Chapter-by-chapter loose-end audit; decision
      register completed through F10; AI-methods appendix (\cref{ch:ml-appendix}). \\
    \bottomrule
  \end{tabular}
\end{table}
